Moderne computertoepassingen bestaan uit onderling verbonden \emph{services} die
worden uitgerold op computers die gedeeld worden door meerdere gebruikers, zoals
in de cloud. Door computers op deze manier te delen, kan het gebruik
geoptimaliseerd en de kosten verlaagd worden. Dit model vereist echter robuuste
beveiligingsoplossingen om vertrouwelijke gegevens te beschermen. Zogeheten
\acp{TEE} beloven een oplossing door gebruik te maken speciale hardware voor
sterke veiligheidsgaranties, zelfs wanneer gevoelige gegevens zich in het
geheugen of de processor bevinden. \acp{TEE} schermen programma's namelijk af
van de rest van het systeem terwijl ze worden uitgevoerd. Zelfs systeemsoftware,
zoals de hypervisor of het besturingssysteem, heeft geen toegang tot de
beveiligde omgeving. Hierdoor worden de mogelijkheden om de beveiliging te
omzeilen aanzienlijk beperkt. Deze aanpak wordt sinds kort omschreven als
\emph{confidential computing}.

In dit proefschrift analyseren we \acp{TEE} vooral vanuit een praktisch
perspectief. We onderzoeken hoe \acp{TEE} gebruikt kunnen worden in zowel
bestaande als nieuwe toepassingen. We bestuderen twee scenario's waarbij
\acp{TEE} kunnen worden ingezet om de beveiliging van het systeem te verbeteren:
smart homes en certificatenbeheer bij \ac{V2X}-toepassingen. Daarnaast bekijken
we of deze technologie, naast veiligheid, ook andere voordelen kan bieden, zoals
verbeterde efficiëntie en schaalbaarheid. We besteden ook aandacht aan de
praktische uitdagingen van deze technologie: we onderzoeken namelijk hoe een
applicatie die binnen een \ac{TEE} draait, interageert met het niet-beschermde
deel van het systeem. We tonen aan hoe ``tijd'' binnen een \ac{TEE}
gemanipuleerd kan worden om de integriteit van applicaties te ondermijnen.
Verder analyseren we ook de vertrouwensrelatie tussen de klant van een
cloud-\ac{TEE} en de cloudprovider, en bespreken we tekortkomingen van het
huidige beveiliginsmodel van \acp{TEE}. Deze technologie heeft het potentieel
een hoeksteen te worden voor de beveiliging van moderne toepassingen. Het
gebruik ervan vereist echter zorgvuldige keuzes van zowel ontwikkelaars als
beheerders, evenals passende ondersteuning van cloudproviders.
